% ============================================================
\chapter{Laws of Large Numbers}
\label{ch:lln}
% ============================================================

Why does the observed frequency of an event stabilize around its theoretical probability when one repeats the experiment many times? This question, at the heart of probabilistic thinking, received its first answer from Jacob Bernoulli in 1713, in his posthumously published \emph{Ars Conjectandi}. Bernoulli showed that the empirical frequency converges to the probability --- but his proof was long and laborious. Chebyshev, in 1867, dramatically simplified the proof using his inequality. Kolmogorov, in 1933, crowned the edifice with the \emph{strong law of large numbers}, which guarantees almost sure convergence. This chapter traces this progression, from concentration inequalities to the various modes of stochastic convergence.

\begin{intuition}
The law of large numbers is one of the most fundamental results in probability
theory. It justifies the frequentist interpretation of probability: the empirical
frequency of an event converges to its theoretical probability as the number of
experiments increases.
\end{intuition}

% ------------------------------------------------------------
\section{Modes of Convergence}
% ------------------------------------------------------------

\begin{definition}[Convergence in probability]
A sequence $(X_n)$ \textbf{converges in probability} to $X$, written
$X_n \xrightarrow{\PP} X$, if for every $\varepsilon > 0$:
\[
    \lim_{n \to \infty} \PP(\abs{X_n - X} > \varepsilon) = 0.
\]
\end{definition}

\begin{definition}[Almost sure convergence]
$(X_n)$ \textbf{converges almost surely} to $X$, written
$X_n \xrightarrow{\text{a.s.}} X$, if
\[
    \PP\!\left(\lim_{n \to \infty} X_n = X\right) = 1,
\]
or equivalently, $\PP\!\left(\bigcap_{m=1}^{\infty}
\bigcup_{n=m}^{\infty}\{\abs{X_n - X} > \varepsilon\}\right) = 0$
for every $\varepsilon > 0$.
\end{definition}

\begin{definition}[Convergence in $L^p$]
$(X_n)$ \textbf{converges in $L^p$} ($p \geq 1$) to $X$, written
$X_n \xrightarrow{L^p} X$, if
\[
    \lim_{n \to \infty} \E[\abs{X_n - X}^p] = 0.
\]
\end{definition}

\begin{theorem}[Hierarchy of convergences]
\[
    X_n \xrightarrow{L^p} X \implies X_n \xrightarrow{\PP} X,
    \qquad
    X_n \xrightarrow{\text{a.s.}} X \implies X_n \xrightarrow{\PP} X.
\]
The converses are false in general. However,
$X_n \xrightarrow{\PP} X$ implies the existence of a subsequence
$X_{n_k} \xrightarrow{\text{a.s.}} X$.
\end{theorem}

\begin{warning}[title=\textbf{a.s.\ convergence and $L^p$ convergence}]
Almost sure convergence does not imply $L^p$ convergence, nor the converse.
Each has its own conditions.
\end{warning}

% ------------------------------------------------------------
\section{Weak Law of Large Numbers}
% ------------------------------------------------------------

\begin{theorem}[Weak law of large numbers (WLLN)]
\label{thm:wlln}
Let $(X_n)_{n \geq 1}$ be a sequence of \textbf{independent and identically
distributed} (i.i.d.) random variables with $\E[X_1] = \mu$ and
$\Var(X_1) = \sigma^2 < \infty$. Set
$\bar{X}_n = \frac{1}{n}\sum_{i=1}^n X_i$. Then
\[
    \bar{X}_n \xrightarrow{\PP} \mu, \quad \text{i.e.} \quad
    \forall\,\varepsilon > 0,\quad
    \lim_{n \to \infty} \PP(\abs{\bar{X}_n - \mu} > \varepsilon) = 0.
\]
\end{theorem}

\begin{proof}
This is a direct application of Chebyshev's inequality.
\begin{align*}
    \PP(\abs{\bar{X}_n - \mu} > \varepsilon)
    &\leq \frac{\Var(\bar{X}_n)}{\varepsilon^2}
    = \frac{\Var\!\left(\frac{1}{n}\sum X_i\right)}{\varepsilon^2}
    = \frac{1}{n^2}\cdot\frac{n\sigma^2}{\varepsilon^2}
    = \frac{\sigma^2}{n\varepsilon^2}
    \xrightarrow[n\to\infty]{} 0. \qedhere
\end{align*}
\end{proof}

\begin{remark}
The WLLN remains true under the weaker assumption that the $X_i$ are uncorrelated
(not necessarily independent) with the same first two moments.
\end{remark}

\begin{intuition}[title=\textbf{Meaning of the WLLN}]
The sample mean $\bar{X}_n$ concentrates around $\mu$ as $n$ grows. This is the
theoretical foundation of Monte Carlo methods and statistical estimation.
\end{intuition}

% ------------------------------------------------------------
\section{Strong Law of Large Numbers}
% ------------------------------------------------------------

\begin{theorem}[Strong law of large numbers (SLLN)]
\label{thm:slln}
Let $(X_n)_{n \geq 1}$ be i.i.d.\ with $\E[\abs{X_1}] < \infty$ and
$\E[X_1] = \mu$. Then
\[
    \bar{X}_n \xrightarrow{\text{a.s.}} \mu, \quad \text{i.e.} \quad
    \PP\!\left(\lim_{n \to \infty} \bar{X}_n = \mu\right) = 1.
\]
\end{theorem}

\begin{remark}
The SLLN requires only the existence of the first moment (no finite variance).
The full proof uses truncation and the Borel--Cantelli lemma.
\end{remark}

\begin{proof}[Proof under the assumption $\E\lbrack X_1^4\rbrack < \infty$]
We give the proof under the stronger assumption $\E[X_1^4] < \infty$.
WLOG assume $\mu = 0$. Set $S_n = \sum_{i=1}^n X_i$.
\begin{align*}
    \E[S_n^4] &= \E\!\left[\left(\sum_{i=1}^n X_i\right)^4\right]
    = \sum_i \E[X_i^4] + \binom{4}{2}\sum_{i \neq j}\E[X_i^2]\E[X_j^2]
\end{align*}
(cross terms with odd powers vanish by independence and $\E[X_i]=0$).
Hence $\E[S_n^4] = n\E[X_1^4] + 3n(n-1)(\E[X_1^2])^2 \leq Cn^2$ for some
constant $C$. By Markov's inequality:
\[
    \PP(\abs{\bar{X}_n} > \varepsilon)
    = \PP(\abs{S_n} > n\varepsilon)
    \leq \frac{\E[S_n^4]}{n^4\varepsilon^4}
    \leq \frac{C}{n^2\varepsilon^4}.
\]
Since $\sum_n 1/n^2 < \infty$, the first Borel--Cantelli lemma gives
$\PP(\limsup\{\abs{\bar{X}_n} > \varepsilon\}) = 0$ for every $\varepsilon > 0$,
so $\bar{X}_n \to 0$ a.s.
\end{proof}

% ------------------------------------------------------------
\section{Applications}
% ------------------------------------------------------------

\subsection{Frequentist interpretation}

\begin{corollary}
Let $A$ be an event with probability $p$. If the experiment is repeated $n$ times
independently, the empirical frequency
$\hat{p}_n = \frac{1}{n}\sum_{i=1}^n \mathbf{1}_{A_i}$ converges almost surely
to $p$.
\end{corollary}

\subsection{Monte Carlo method}

\begin{example}[Estimating $\pi$]
Let $(X_i, Y_i)_{i \geq 1}$ be i.i.d.\ uniform on $[0,1]^2$. Set
$Z_i = \mathbf{1}_{X_i^2+Y_i^2 \leq 1}$. Then $\E[Z_i] = \pi/4$ and by the
law of large numbers:
\[
    \frac{4}{n}\sum_{i=1}^n Z_i \xrightarrow{\text{a.s.}} \pi.
\]
\end{example}

\subsection{Glivenko--Cantelli theorem}

\begin{theorem}[Glivenko--Cantelli]
Let $(X_i)$ be i.i.d.\ with CDF $F$ and let
$\hat{F}_n(x) = \frac{1}{n}\sum_{i=1}^n \mathbf{1}_{X_i \leq x}$ be the
empirical CDF. Then
\[
    \sup_{x \in \R}\abs{\hat{F}_n(x) - F(x)} \xrightarrow{\text{a.s.}} 0.
\]
\end{theorem}

% ------------------------------------------------------------
\section{Convergence in Distribution (Introduction)}
% ------------------------------------------------------------

\begin{definition}[Convergence in distribution]
$(X_n)$ \textbf{converges in distribution} to $X$, written
$X_n \xrightarrow{\mathcal{L}} X$, if for every bounded continuous
$f : \R \to \R$:
\[
    \lim_{n \to \infty}\E[f(X_n)] = \E[f(X)].
\]
Equivalently: $F_{X_n}(x) \to F_X(x)$ at every continuity point $x$ of $F_X$.
\end{definition}

\begin{theorem}[Complete hierarchy]
\[
    X_n \xrightarrow{\text{a.s.}} X \implies X_n \xrightarrow{\PP} X
    \implies X_n \xrightarrow{\mathcal{L}} X.
\]
Convergence in distribution is the weakest. It concerns only the laws, not the
variables themselves.
\end{theorem}

\begin{keyformulas}[title=\textbf{Summary of convergences}]
\begin{center}
\begin{tabular}{ll}
    \toprule
    \textbf{Convergence} & \textbf{Definition} \\
    \midrule
    Almost sure & $\PP(\lim X_n = X) = 1$ \\
    In probability & $\PP(\abs{X_n - X} > \varepsilon) \to 0$ \\
    In $L^p$ & $\E[\abs{X_n - X}^p] \to 0$ \\
    In distribution & $\E[f(X_n)] \to \E[f(X)]$ for bounded continuous $f$ \\
    \bottomrule
\end{tabular}
\end{center}
\end{keyformulas}

% ------------------------------------------------------------
\section{Exercises}
% ------------------------------------------------------------

\begin{exercise}
Let $X_n \sim \mathcal{U}([0, 1/n])$. Show that $X_n \xrightarrow{\PP} 0$
and $X_n \xrightarrow{L^2} 0$. Is the convergence almost sure?
\end{exercise}

\begin{exercise}
Let $(X_n)$ be i.i.d.\ with $\PP(X_n = 1) = \PP(X_n = -1) = 1/2$.
Show that $\bar{X}_n \to 0$ a.s.\ and that $\PP(\bar{X}_n = 0) \to 0$ when $n$
is odd. Comment.
\end{exercise}

\begin{exercise}
Give an example of a sequence $(X_n)$ converging in probability to $0$ but not
almost surely.
\textit{Hint:} consider indicators of intervals sliding over $[0,1]$.
\end{exercise}

\begin{exercise}
Draw $n$ numbers independently and uniformly from $[0,1]$. Let
$M_n = \max(X_1, \ldots, X_n)$. Show that $M_n \to 1$ a.s.
\end{exercise}

\begin{exercise}
Let $(X_n)$ be i.i.d.\ $\mathrm{Exp}(1)$. Show that
$\frac{1}{n}\max(X_1, \ldots, X_n) \to 0$ a.s., and even that
$\max(X_1,\ldots,X_n)/\ln n \to 1$ a.s.
\end{exercise}

\begin{exercise}
Let $(X_n)$ be i.i.d.\ standard Cauchy. Show that $\bar{X}_n$ does not converge
in probability to any constant.
\textit{Hint:} use characteristic functions.
\end{exercise}

\begin{exercise}[Monte Carlo]
We wish to estimate $I = \int_0^1 e^{-x^2}\,dx$ by Monte Carlo. Write the
Monte Carlo estimator and, using Chebyshev's inequality, determine the number of
samples needed so that the error is less than $0.01$ with probability $0.95$.
\end{exercise}

\begin{exercise}
Let $(X_n)$ be i.i.d.\ with $\E[X_1] = 0$, $\Var(X_1) = 1$. Show by the WLLN
that $S_n/n \to 0$ in probability, then that $S_n^2/n^2 \to 0$ in probability.
Deduce a result about $S_n/n^{3/4}$.
\end{exercise}
